
%%%%%%%%%%%%%%%%%%%%%%%%%%%%%%%%%%%%%%%%%%%%%%%%%%%%%%%
%%%%%%%%%%%%%%%%%%%%%%%%%%%%%%%%%%%%%%%%%%%%%%%%%%%%%%%
%%%%%%%%%%%%%%%%%%%%%%%%%%%%%%%%%%%%%%%%%%%%%%%%%%%%%%%





%%%%%%%%%%%%%%%%%%%%%%%%%%%%%%%%%%%%%%%%%%%%%%%%%%%%%%%
%%%%%%%%%%%%%%%%%%%%%%%%%%%%%%%%%%%%%%%%%%%%%%%%%%%%%%%
%%%%%%%%%%%%%%%%%%%%%%%%%%%%%%%%%%%%%%%%%%%%%%%%%%%%%%%

\begin{frame}[t]{Generating $\ytil$}
%
\begin{itemize}
\item We have defined balance checks in terms of $\ytil$ such that
$\expect{}{\ytil \vert \x} = \expect{}{\y \vert \x} + \delta f(\x)$
\item We have approximated $\muhatmrp[\Ytil] - \muhatmrp$ for $\ytil \approx \y$
\end{itemize}
%

How to get such a $\ytil$?  \textbf{Recall $\y$ is binary!}  \onslide<2->{
    \textbf{Two solutions, with their own pros and cons:}
}
\splitpagenoline{
    \centering
    \onslide<2->{
        \textbf{Option 1:} Force $\ytil$ to be binary.
    }
    %
    \onslide<3->{
        \begin{enumerate}
        \item Make \textit{some} guess $\hat{m}(\x) \approx \expect{}{\y \vert \x}$
        %
        \begin{itemize}
        \item E.g.~Posterior mean, or
        \item Shrunken posterior mean, or
        \item Some values that gives the same posterior
        \end{itemize}
        %
        \item Take \surcol{$u_i \iid \mathrm{Unif}(0,1)$}
        \item Assume \surcol{$\y_i = \ind{u_i \le \hat{m}(\x_i)}$}
        \item Draw \surcol{$u_n \vert \y_n$}
        \item Set \surcol{$\ytil_i = \ind{u_i \le \hat{m}(\x_i)+ \delta \x_i}$}
    }
%
\end{enumerate}
%
} {
    \centering
    \onslide<2->{
        \textbf{Option 2:} Allow $\ytil$ to take generic values.
    }
    \onslide<4->{
        \begin{enumerate}
            \item Set \surcol{$\ytil_i = \y_i + \delta \f(\x_i)$}.
            \item Then you're done.
            \item There is nothing else to do.
            \item This space deliberately left blank.
        \end{enumerate}
    }
}
%
\\[1em]
\onslide<5->{
    \splitpagenoline{
    %
\textbf{Pros and cons:}
    \begin{itemize}
    \item Realistic
    \item Have to pick $\hat{m}(\x)$
    \item \surcol{$\Ytil - \Ysur$} not infinitesimally small
    \item \textbf{Use for checks \& experiments}
    \end{itemize}
    %
    }
    {
\textbf{Pros and cons:}
    \begin{itemize}
    \item Not realistic
    \item No additional assumptions
    \item \surcol{$\Ytil - \Ysur$} may be infinitesimally small
    \item \textbf{Use for theory}
    \end{itemize}
    }
}



\end{frame}



%%%%%%%%%%%%%%%%%%%%%%%%%%%%%%%%%%%%%%%%%%%%%%%%%%%%%%%
%%%%%%%%%%%%%%%%%%%%%%%%%%%%%%%%%%%%%%%%%%%%%%%%%%%%%%%
%%%%%%%%%%%%%%%%%%%%%%%%%%%%%%%%%%%%%%%%%%%%%%%%%%%%%%%

\begin{frame}[c]{Theory}

\textbf{When is the local approximation accurate?}

\begin{block}{Theorem \parencite{giordano:2026:mrplew}:}
    Take $\surcol{\ytil_i = \y_i + \deltaf \f(\x_i)}$.\\[1em]
    We state conditions for Bayesian hierarchical logistic regression under which
$$
\onslide<4->{\sup_{f \in \mathcal{F}}}
\abs{
    \muhatmrp[\Ytil] - \muhatmrp
    - \deltaf \surcol{\sumsur \w_i^\mrp \f(\x_i)}
}
\only<1>{ = \textrm{Small}}
\only<2->{ = O(\deltaf^2)}
\onslide<3->{\textrm{ as }N \rightarrow \infty}
$$
\onslide<4->{
    ...for a very broad class of $\mathcal{F}$.
    \footnote{$\mathcal{F}$ can be any Donsker class of measurable functions
    with uniformly bounded
    $\expect{}{\norm{\x} \f(\x)}_2^2 $.}
}
\end{block}

\onslide<4->{
    \textbf{Uniformity justifies searching for ``imbalanced'' $\f$.}
}

\onslide<5->{
The uniformity result builds on our earlier work on uniform
and finite--sample error bounds
for Bernstein--von Mises theorem--like results
\footcite{giordano:2024:bayesij,kasprzak:2025:laplace}.
}

\end{frame}


%%%%%%%%%%%%%%%%%%%%%%%%%%%%%%%%%%%%%%%%%%%%%%%%%%%%%%%%%%%%%%%%%
%%%%%%%%%%%%%%%%%%%%%%%%%%%%%%%%%%%%%%%%%%%%%%%%%%%%%%%%%%%%%%%%%
%%%%%%%%%%%%%%%%%%%%%%%%%%%%%%%%%%%%%%%%%%%%%%%%%%%%%%%%%%%%%%%%%





\begin{frame}[t]{Predictions and actual MCMC results}
    \AlexanderRefitPlot{}
\end{frame}




\begin{frame}[t]{Predictions and actual MCMC results}
    \LaxRefitPlot{}

    \begin{center}
    More careful investigation shows that the linear
    approximation is computed correctly, but the posterior
    is highly curved in this direciton, even for small
    step sizes.\\[1em]
    \textbf{There is real work to do to understand this!}
    \end{center}
\end{frame}








%%%%%%%%%%%%%%%%%%%%%%%%%%%%%%%%%%%%%%%%%%%%%%%%%%%%%%%%%%%%%%%%%
%%%%%%%%%%%%%%%%%%%%%%%%%%%%%%%%%%%%%%%%%%%%%%%%%%%%%%%%%%%%%%%%%
%%%%%%%%%%%%%%%%%%%%%%%%%%%%%%%%%%%%%%%%%%%%%%%%%%%%%%%%%%%%%%%%%

% \begin{frame}[c]{Other uses}
% \centering


% Applying the idea to the observations negative weights measures
% \emph{non-monotonicity}.

% \wholeslidefig{
% \AlexanderWeightPlot{}
% }

% % \splitpagenoline{
% %     \AlexanderWeightPlot{}
% % }{
% %     \LaxWeightPlot{}
% % }
% \end{frame}
